\section{内核级扫描结果明细}
\label{app:conescan}

\subsection{导出文件概况（conescan 第一遍）}

\begin{center}\small
\begin{tabular}{@{}lr@{\qquad}lr@{}}
\toprule
\multicolumn{2}{c}{\textbf{总量}} & \multicolumn{2}{c}{\textbf{声明记录}}\\
\midrule
行数 & 1,080,769,508 & \lean{thm} & 1,246,110 \\
字节数 & 61,570,595,788 & \lean{def} & 214,241 \\
名字 & 14,065,150 & \lean{inductive}（声明块） & 7,969 \\
表达式结点 & 1,065,214,335 & \lean{opaque} & 91 \\
其中 \lean{app} & 989,707,003 & \lean{quot} & 4 \\
其中 \lean{proj} & 65,718 & \lean{axiom} & 3 \\
标记 unsafe／partial 的记录 & 0 & \textbf{合计} & \textbf{1,468,418} \\
\bottomrule
\end{tabular}
\end{center}

公理：\lean{propext}、\lean{Quot.sound}、\lean{Classical.choice}。导出器：lean4export 3.1.0；Lean 4.33.1（githash \texttt{819816b2}）。

\subsection{依赖 \texttt{anyGeometry} 的全部 30 个声明}

「类型／值」两列表示 \lean{anyGeometry}（经传递依赖）出现在该声明的类型里还是证明项（值）里；「直接见证」是该声明直接引用的一个已被污染的常量。

{\scriptsize
\begin{longtable}{@{}rlcc>{\raggedright\arraybackslash}p{6.9cm}>{\raggedright\arraybackslash}p{5.4cm}@{}}
\caption{内核级：传递依赖 \texttt{GlobalTheory.Mem.Top.Bl.anyGeometry} 的声明}\label{tab:tainted}\\
\toprule
\# & 种类 & 类型 & 值 & 声明 & 直接见证 \\
\midrule
\endfirsthead
\toprule
\# & 种类 & 类型 & 值 & 声明 & 直接见证 \\
\midrule
\endhead
\bottomrule
\endfoot
\tableinput{tables/tainted_anygeometry.tex}
\end{longtable}
}

\subsection{读取 \texttt{OAt} 的 \texttt{geo} 字段的声明}

直接读取该字段的只有两个定理：\lean{divEx_geoCl} 与 \lean{oAt_setDiv}（表中「直接见证」为 \lean{OAt.geo} 的两行）；其余各行是它们的传递使用者。

{\scriptsize
\begin{longtable}{@{}rlcc>{\raggedright\arraybackslash}p{6.9cm}>{\raggedright\arraybackslash}p{5.4cm}@{}}
\caption{内核级：传递依赖「对 \texttt{OAt} 第 0 字段（\texttt{geo}）的投影」的声明}\\
\toprule
\# & 种类 & 类型 & 值 & 声明 & 直接见证 \\
\midrule
\endfirsthead
\toprule
\# & 种类 & 类型 & 值 & 声明 & 直接见证 \\
\midrule
\endhead
\bottomrule
\endfoot
\tableinput{tables/tainted_oatgeo.tex}
\end{longtable}
}

\subsection{引用 \texttt{GeoCl.rec} 的声明}

只有 \lean{GeoCl} 自身的归纳声明块（类型、两个构造子与递归子共 4 条记录，递归子的类型提到它自己）。没有任何其他声明。

\scanThreeAppendix

\section{其余源码摘录}

从 \path{Main.lean} 到「文献定理」层的三级转接：

\excerpt{maintheorems}
\excerpt{mainconcl}
\excerpt{lit810}

\section{完整哈希与版本}

\begin{center}\scriptsize
\renewcommand{\arraystretch}{1.25}
\begin{tabularx}{\textwidth}{@{}>{\raggedright\arraybackslash}p{3.6cm}>{\raggedright\arraybackslash}X@{}}
\toprule
对象 & SHA-256／提交 \\
\midrule
\path{lean.tar.gz}（作者原始包） & \texttt{9631b4f020199e580ba86178cb56c6e9b09f690c2369a80dc4334aa93f49c7fc} \\
\path{solution.ndjson.zst} & \texttt{61a8938ad5f5e57397c1c0e3b954916af546d9004ade01c3555744a681b05381} \\
\path{solution.ndjson} & \texttt{5ea74b912147c5c93cc2b6d70cc727cf908df2f62ac6a20480a3fa02f6be244e} \\
\path{challenge.ndjson}（公布的与本机重建的） & \texttt{bd484ac9886a96f26b1d9a7a59ebd45b7a70d102c015d4f31f9302bbdd6ab2b9} \\
镜像 \path{the-omega-institute/riemann-hypothesis} & 提交 \texttt{c7d25b3bfd7906388bb8eb0774e027a2b756a7b6} \\
Mathlib & \texttt{0df444a360eaa60ab8c11dca51a86af692955474}（官方标签 \texttt{v4.33.1}） \\
con-leche & \texttt{bc6676101c488e21ce71a317b33905534d497e14}，以 Lean \texttt{v4.36.0}（\texttt{611c117b}）编译 \\
lean4export & \texttt{15f6055e299ad5b89345e533cc2192f4cc00f659}（标签 \texttt{v4.33.0}），以 Lean \texttt{v4.33.1}（\texttt{819816b2}）编译 \\
\bottomrule
\end{tabularx}
\end{center}

\section{复现步骤}

仓库 \path{evidence/} 下的 \path{REPRODUCE.md} 给出了逐条命令。概要如下（需要约 80\,GB 磁盘；con-leche 一步在 16 核机器上约需一小时）：

\begin{Verbatim}[fontsize=\scriptsize,frame=single,rulecolor=\color{gray!50}]
# 1. 源码：下载 alpo.ge 的 14 个分卷，拼接并核对 SHA-256，解包
for i in 01 02 ... 14; do curl -LO https://alpo.ge/lean.tar.gz.part$i; done
cat lean.tar.gz.part?? > lean.tar.gz && shasum -a 256 lean.tar.gz

# 2. 源码层面：定位定义与调用链
rg -n -w anyGeometry src | wc -l          # 118 行，24 个文件
rg -n 'everyE₂|plugAll' src               # plugAll 只被 everyE₂ 调用；everyE₂ 无调用者

# 3. 内核层面：下载导出文件，核对哈希，运行 conescan
curl -LO https://math.mit.edu/~drew/grh/solution.ndjson.zst
zstd -dc --long=30 solution.ndjson.zst | conescan OUTDIR seeds.tsv watch.txt

# 4. 独立内核重放与命题比对
con-leche --verified --compare challenge.ndjson solution.ndjson

# 5. 独立重建挑战导出（官方 Lean v4.33.1 + Mathlib 官方缓存 + lean4export v4.33.0）
lake exe cache get && lake build Challenge
lake env lean4export Challenge -- $(cat export-targets.txt) | shasum -a 256
\end{Verbatim}

\section{参考来源}

{\small\sloppy
\begin{itemize}
\item L.~Alpöge, 公告：\url{https://alpo.ge/remark.pdf}；源码分卷：\url{https://alpo.ge/lean.tar.gz.part01} 至 \texttt{part14}。
\item J.~L.~Nielsen, \emph{Invalidity of the Alpoge Proposed Proof of the Riemann Hypothesis}：\url{https://philarchive.org/rec/NIEIOT-2}（本次审计未能取得全文）；
作者在 X 上的表述：\url{https://x.com/QualiaQuanta/status/2123871320949825776}。
\item frankchu, \emph{GRH proof's geometry is \texttt{fun \_ \_ \_ => True}, in 24 files}, DEV Community：\url{https://dev.to/frankchu/grh-proofs-geometry-is-fun-true-in-24-files-2n2b}。
\item Lean 与 Mathlib 社区联合声明：\url{https://leanprover-community.github.io/blog/posts/riemann-hypothesis/}。
\item A.~Sutherland, \emph{Verification files for the Lean proof of the Riemann Hypothesis}：\url{https://math.mit.edu/~drew/grh/}。
\item 源码镜像：\url{https://github.com/the-omega-institute/riemann-hypothesis}；社区验证仓库：\url{https://github.com/Deicyde/riemann-hypothesis}。
\item 检查器：\url{https://github.com/leanprover/con-leche}；导出器：\url{https://github.com/leanprover/lean4export}。
\end{itemize}
}
